\section{Completion of the paving argument}\label{sec:completion}

We now combine the coloring theorem and moment transfer. The coloring estimates apply to finite-support approximations of a Cartan operator. Before transferring a fixed moment, we must return to the original operator without requiring a rate of approximation depending exponentially on the moment order. A clipping argument accomplishes this. Its use in a nontracial state requires the following elementary observation about varying right multipliers.

\subsection{Continuous calculus for varying sequences}

\begin{lemma}\label{lem:multiplier-calculus}
Let $\mathcal L$ be a von Neumann algebra with a faithful normal state $\omega$. Let $\mathscr R_\omega$ consist of all uniformly bounded sequences $(R_n)$ in $\mathcal L$ such that
\[
 \|D_nR_n\|_\omega\longrightarrow0
\]
whenever $(D_n)$ is uniformly bounded and $\|D_n\|_\omega\to0$. Then $\mathscr R_\omega$ is an algebra, closed for the supremum operator norm. It contains all constant sequences and all uniformly bounded sequences in the centralizer of $\omega$. It is also closed under adding uniformly bounded sequences converging to zero in $\|\cdot\|_\omega$.

If $(U_n),(V_n)\in\mathscr R_\omega$ are self-adjoint, uniformly bounded, and $\|U_n-V_n\|_\omega\to0$, then for every continuous function $f$ on a compact interval containing their spectra,
\[
 (f(U_n)),(f(V_n))\in\mathscr R_\omega,
 \qquad \|f(U_n)-f(V_n)\|_\omega\longrightarrow0.
\]
\end{lemma}

\begin{proof}
The assertions for constant sequences and centralizer sequences follow from Lemma~\ref{lem:state-multipliers}. If $(R_n),(S_n)\in\mathscr R_\omega$ and $(D_n)$ is a bounded state-norm-null sequence, then $(D_nR_n)$ is another such sequence, so $\|D_nR_nS_n\|_\omega\to0$. Addition is immediate. Closure in the supremum norm follows by approximating the right multiplier uniformly before taking the limit in $n$.

For the perturbation assertion, if $(E_n)$ is bounded and state-norm-null, then
\[
 \|D_nE_n\|_\omega\le\|D_n\|\,\|E_n\|_\omega\longrightarrow0.
\]
Thus $(R_n+E_n)\in\mathscr R_\omega$ whenever $(R_n)\in\mathscr R_\omega$.

For a polynomial $f$, the assertion about membership follows from the algebra property. Polynomial differences are controlled by the telescoping identity
\[
 U_n^j-V_n^j=\sum_{h=0}^{j-1}U_n^h(U_n-V_n)V_n^{j-1-h}.
\]
Every right multiplier on the right belongs to $\mathscr R_\omega$, while the left multipliers are uniformly bounded. Each summand therefore tends to zero in state norm. Uniform polynomial approximation on the common spectral interval proves both assertions for continuous $f$.
\end{proof}

\subsection{Paving in the presence of a central state}

\begin{proposition}\label{prop:expected}
Let $A\subseteq M$ be a masa and let $\varphi$ be a faithful normal state on $M$ with $A$ in its centralizer. For $0<\eps<1$, put
\[
 r_{\mathrm e}=\left\lceil 20000^2\eps^{-2}\right\rceil.
\]
For every self-adjoint contraction $x\in M$ and every $\rho>0$, there are a partition $P=(p_1,\ldots,p_{r_{\mathrm e}})$ in $A$ and a projection $q\in M$ such that
\[
 \varphi(1-q)<\rho,
 \qquad \|q(S_P(x)-E_Ax)q\|\le\eps.
\]
Here $E_A$ is the $\varphi$-preserving conditional expectation onto $A$.
\end{proposition}

\begin{proof}
By Lemma~\ref{lem:separable-reduction}, it suffices to work with separable predual. The conditional expectation in that reduction is the restriction of $E_A$, so proving the stated diagonal choice in the smaller inclusion proves it in the original one. We henceforth use the notation of Sections~\ref{sec:kernels} and~\ref{sec:transfer}.

Fix $r=r_{\mathrm e}$ and write
\[
 a=E_Ax,\qquad y=E_Nx-E_Ax\in N,
 \qquad w=x-E_Nx\in\ker E_N.
\]
The three elements are self-adjoint, $\|a\|\le1$, and $\|y\|,\|w\|\le2$. Moreover $E_Ay=0$: uniqueness of the $\varphi$-preserving expectation onto $A$ gives $E_AE_N=E_A$.

\paragraph{Finite-support Cartan approximations.}
Choose self-adjoint contractions in $N_{\mathrm{alg}}$ converging strongly to $E_Nx$. To obtain them, first apply Kaplansky density to the norm closure of the finite normalizer algebra, then approximate its self-adjoint contractions in norm by self-adjoint elements of $N_{\mathrm{alg}}$ and rescale. Bounded strong approximation can be taken sequentially in the faithful separable state representation.

For each approximant, only finitely many partial isomorphisms occur in its support. We may also arrange a bound on their edge Radon--Nikodym ratios. Indeed, for such a finite list of maps $\gamma$, remove from $X$ each domain set where $\Delta(\gamma v,v)$ lies outside $[b^{-1},b]$, and let $e_b\in A$ be the indicator of the remaining set. As $b\to\infty$, these projections increase strongly to $1$. Compressing the approximant by $e_b$ preserves self-adjointness and its contraction bound, and restricts every surviving edge to the required ratio range. Enlarge the finite list by its inverses if necessary. Choosing $b$ sufficiently large for each approximant preserves strong convergence.

Subtracting the diagonal now gives a sequence
\begin{equation}\label{eq:cartan-approximants}
 y_n=y_n^*\in N_{\mathrm{alg}},\qquad E_Ay_n=0,
 \qquad \|y_n\|\le2,\qquad \|y_n-y\|_\varphi\longrightarrow0.
\end{equation}
Each $y_n/2$ is represented by measurable zero-diagonal Hermitian contractions on orbit $\ell^2$ spaces, supported on an undirected graph of finite degree with bounded edge cocycle ratios. These bounds may depend on $n$. Theorem~\ref{thm:bernoulli-coloring} permits this dependence: it does not affect $r$.

\paragraph{One partition for the two operator terms.}
Use the fixed relative Bernoulli extension $\widetilde N$ and amalgamated product
\[
 \mathcal K=M*_N\widetilde N,
 \qquad \Phi=\varphi|_N\circ E_N^{\mathcal K},
\]
from Lemma~\ref{lem:bernoulli-model}. Apply Theorem~\ref{thm:bernoulli-coloring} to $y_n/2$ with
\begin{equation}\label{eq:coloring-parameters}
 t=\frac4{\sqrt r},\qquad \alpha=\frac\eps{32},
 \qquad m=n,\qquad \eta=\left(\frac\eps{16}\right)^{2n}.
\end{equation}
Denote the color projections by $p_1^n,\ldots,p_r^n\in B$. They form a partition, belong to the centralizer of $\Phi$, and satisfy
\begin{equation}\label{eq:color-marginals-completion}
 \widetilde E_N(p_i^n)=r^{-1}1\qquad(1\le i\le r).
\end{equation}
With $\beta=80$, the norm bound supplied by the coloring theorem, after restoring the factor two, obeys
\begin{equation}\label{eq:cartan-numerical-bound}
 2(2\beta t+\alpha)=\frac{1280}{\sqrt r}+\frac\eps{16}
 \le\frac{253}{2000}\eps<\frac\eps4.
\end{equation}

Put
\[
 \widehat Y_n=\sum_{i=1}^r p_i^n y_n p_i^n\in\widetilde N.
\]
At a root where the radius-$n$ ball has the asserted pinching bound, $\widehat Y_n^n$ applied to the root vector agrees with the $n$th power of that principal matrix applied to the same vector. Every path of length $n$ starting at the root stays in its radius-$n$ ball. At all other roots, $\|\widehat Y_n\|\le2$. Integrating the squared norms of these vectors gives
\begin{equation}\label{eq:cartan-high-moment}
 \Phi(\widehat Y_n^{2n})
 \le\left(\frac\eps4\right)^{2n}
       +2^{2n}\left(\frac\eps{16}\right)^{2n}.
\end{equation}
This uses diagonal integration in the relation representation, not a trace on the nonsingular relation algebra.

For the original operators set
\[
 Y_n=\sum_{i=1}^r p_i^n y p_i^n,
 \qquad W_n=\sum_{i=1}^r p_i^n w p_i^n.
\]
Pinching by a centralizer partition is contractive in the state norm, so Equation~\eqref{eq:cartan-approximants} implies
\begin{equation}\label{eq:pinched-approximation}
 \|Y_n-\widehat Y_n\|_\Phi\longrightarrow0.
\end{equation}
By Lemma~\ref{lem:free-compression} and Equation~\eqref{eq:color-marginals-completion},
\begin{equation}\label{eq:free-numerical-bound}
 \|W_n\|\le2\left(\frac1r+\frac2{\sqrt r}\right)
 \le\frac{40001}{200000000}\eps<\frac\eps4.
\end{equation}
Here orthogonality of the $p_i^n$ makes the norm of the block sum the maximum of its block norms. Thus the same $r$ projections control the complementary term; there is no second partition refinement.

\paragraph{Clipping before fixing a transferred moment.}
Let $f:\mathbb R\to[-\eps/2,\eps/2]$ be the continuous clipping function
\[
 f(s)=\max\{-\eps/2,\min\{s,\eps/2\}\}.
\]
Spectral calculus and Equation~\eqref{eq:cartan-high-moment} show that
\[
 \Phi\bigl(1_{\{|\widehat Y_n|>\eps/2\}}\bigr)
 \le\left(\frac2\eps\right)^{2n}\Phi(\widehat Y_n^{2n})
 \le2^{-2n}+4^{-2n}.
\]
Since $\|\widehat Y_n\|\le2$, it follows that
\begin{equation}\label{eq:approximate-clipping}
 \|\widehat Y_n-f(\widehat Y_n)\|_\Phi\longrightarrow0.
\end{equation}

We apply Lemma~\ref{lem:multiplier-calculus} in $\mathcal K$. Each sequence $(p_i^n)$ lies in $\mathscr R_\Phi$, and $r$ is fixed. Hence $(Y_n)$ and $(W_n)$ belong to that algebra, being finite sums of products of centralizer sequences and fixed operators. Equation~\eqref{eq:pinched-approximation} also places $(\widehat Y_n)$ in it. Equations~\eqref{eq:pinched-approximation}--\eqref{eq:approximate-clipping} and the continuous-calculus conclusion of the lemma yield
\begin{equation}\label{eq:original-clipping}
 \|Y_n-f(Y_n)\|_\Phi\longrightarrow0.
\end{equation}

The self-adjoint operators
\[
 Z_n=Y_n+W_n=\sum_{i=1}^r p_i^n(x-a)p_i^n
\]
therefore differ by a bounded state-norm-null sequence from
\[
 Q_n=f(Y_n)+W_n,\qquad \|Q_n\|\le\frac{3\eps}{4}.
\]
Both sequences belong to $\mathscr R_\Phi$. For every fixed integer $k\ge1$, telescoping their powers gives
\begin{equation}\label{eq:fixed-model-moment}
 \limsup_{n\to\infty}\Phi(Z_n^{2k})
 \le\left(\frac{3\eps}{4}\right)^{2k}.
\end{equation}
Only a fixed power is used at this stage. In particular, Equation~\eqref{eq:cartan-approximants} required no prescribed rate of convergence.

\paragraph{Transfer and the final spectral cut.}
Fix $\rho>0$. First choose $k$ so large that $(3/4)^{2k}<\rho/4$. Then choose $n$ sufficiently large in Equation~\eqref{eq:fixed-model-moment} that
\[
 \Phi(Z_n^{2k})<\frac\rho2\eps^{2k}.
\]
With $n$ and $k$ fixed, approximate the partition $(p_i^n)$ in $B$ by a finite-symbol cylinder partition, using Lemma~\ref{lem:cylinder-partitions}. Bounded state-norm approximation of these self-adjoint projections is strong-* approximation, so the fixed polynomial moment changes arbitrarily little. We may retain the strict upper bound $\rho\eps^{2k}$ with positive slack.

Proposition~\ref{prop:moment-transfer} applies to the expansion of this moment. For a sufficiently fine base partition, its averaged evaluated moment remains strictly less than $\rho\eps^{2k}$. Every evaluation is a genuine projection partition in $A$, and its even moment is nonnegative. Some deterministic evaluation consequently gives a partition $P=(p_1,\ldots,p_r)$ with
\[
 \varphi\left(\left[\sum_{i=1}^r p_i(x-a)p_i\right]^{2k}\right)
 <\rho\eps^{2k}.
\]
Let $z=\sum_i p_i(x-a)p_i=z^*$ and take $q=1_{[-\eps,\eps]}(z)\in M$. Then
\[
 \varphi(1-q)\le\eps^{-2k}\varphi(z^{2k})<\rho,
 \qquad \|qzq\|\le\eps.
\]
Because $a\in A$ commutes with the partition, $z=S_P(x)-a$. This proves the proposition with the number of colors fixed before $\rho$.
\end{proof}

\subsection{Arbitrary inclusions}

\begin{proof}[Proof of Theorem~\ref{thm:main}]
If $x=0$, take $p_1=1$, $p_2=\cdots=p_{r_\eps}=0$, $a=0$, and $q=1$. Otherwise replace $x$ by $x/\|x\|$. Proposition~\ref{prop:expected}, together with Proposition~\ref{prop:state-reduction}, supplies the required paving in the original representation using
\[
 r=\left\lceil400000000\eps^{-2}\right\rceil
      +\left\lceil4/\eps\right\rceil+1
\]
projections. This number depends only on $\eps$. The supporting corner and all partitions, diagonal elements, and compression projections may depend on $F$ and $\delta$, as permitted in Definition~\ref{def:so-paving}. The construction has diagonal norm at most one, and
\[
 r\le400000000\eps^{-2}+4\eps^{-1}+3
 \le400000007\eps^{-2}<500000000\eps^{-2}.
\]
Multiplying the diagonal by the original $\|x\|$ restores scale and gives every assertion of Theorem~\ref{thm:main}.
\end{proof}
